% ----------------------------------------------------------
\chapter{Introdução}
% ----------------------------------------------------------

O concreto armado (CA) é um material fundamental na engenharia civil, com aplicação em uma vasta gama de construções, desde edifícios residenciais e comerciais até obras de infraestrutura de grande porte, como pontes e viadutos. Sua combinação da resistência à compressão do concreto com a ductilidade do aço confere aos elementos estruturais alta capacidade resistente e versatilidade construtiva. Entretanto, essas estruturas permanecem constantemente expostas a ambientes agressivos, tornando-as suscetíveis a processos de degradação que comprometem seu desempenho ao longo do tempo. Entre os mais relevantes, destaca-se a corrosão, que afeta diretamente a integridade das armaduras internas. 

A exposição contínua a esse mecanismo pode resultar em fissuração, perda de seção da armadura, destacamento do cobrimento e, consequentemente, redução da capacidade resistente de estruturas de concreto armado \cite{vu_structural_2000, pereira_junior_reliability_2023, soltani_state---art_2019, possan_co2_2016}. A perda de seção das armaduras longitudinais resulta na redução da rigidez e da resistência à flexão, enquanto a deterioração das armaduras transversais compromete a capacidade de cisalhamento e altera os modos de falha das vigas. Além disso, a degradação da aderência entre o concreto e o aço, induzida pela corrosão, prejudica a transferência de esforços internos, contribuindo para um desempenho estrutural globalmente inferior \cite{soltani_state---art_2019}.

A corrosão das armaduras em vigas de concreto armado pode ser classificada em dois tipos principais: a corrosão induzida por carbonatação e a induzida por cloretos. A carbonatação é um processo químico no qual o dióxido de carbono (CO\(_2\)) presente na atmosfera penetra nos poros do concreto e reage com o hidróxido de cálcio (Ca(OH)\(_2\)), formando carbonato de cálcio (CaCO\(_3\)) e reduzindo o pH da matriz cimentícia, conforme descrito por \textcite{helene_p_r_l_concreto_1993, ortega_assessment_2016}. Essa queda de pH compromete a camada passiva que protege as armaduras, dando início ao processo corrosivo \cite{possan_co2_2016}. Já a penetração de cloretos, mais comum em ambientes marinhos ou regiões onde se utilizam sais descongelantes, acelera a corrosão por meio de mecanismos de difusão, conforme descrito por \textcite{vu_structural_2000}, com base na Lei de Fick.

Nesse cenário, a corrosão impacta diretamente a vida útil remanescente (VUR) das estruturas, tornando a previsão desse parâmetro uma tarefa essencial para assegurar a segurança estrutural e planejar intervenções de manutenção de forma eficiente. A VUR tem como objetivo estimar o intervalo de tempo residual no qual uma estrutura ainda poderá desempenhar sua função com segurança. Essa estimativa parte do estado atual da estrutura e se estende até o momento em que ela perde os requisitos que definem sua durabilidade: a segurança e a funcionalidade. Segundo \textcite{bertolini_corrosion_2013} e \textcite{mehta_concrete_2014}, a durabilidade de uma estrutura está diretamente relacionada à capacidade do concreto em proteger as armaduras contra a corrosão por carbonatação e em manter suas propriedades mecânicas sob exposição a ambientes agressivos.

Na prática, a avaliação da VUR permite transformar dados de degradação, como profundidade de carbonatação, perda de seção das armaduras ou aumento da probabilidade de falha, em estimativas quantitativas do tempo restante até a necessidade de intervenção. \textcite{bertolini_corrosion_2013, okoh_overview_2014} destacam que modelos baseados em confiabilidade, que combinam a evolução da carbonatação com a perda de capacidade resistente e a redução do índice de confiabilidade estrutural, são ferramentas eficazes para essa estimativa. Além disso, estudos recentes \cite{silva_statistical_2016, zhang_prediction_2018} apontam que a integração entre monitoramento, simulação probabilística e técnicas de amostragem, como Monte Carlo, fornece uma base metodológica robusta para decisões de manutenção baseadas em risco. 

Nesse contexto, a previsão da vida útil remanescente constitui o principal objetivo deste trabalho, por meio do desenvolvimento de um modelo preditivo capaz de estimar o tempo em que o índice de confiabilidade atinge o limiar mínimo aceitável para estruturas de concreto armado. O modelo será capaz de considerar o tempo já decorrido de utilização da estrutura e, a partir disso, estimar o período de vida útil restante. As metodologias probabilísticas destacam-se nesse processo por oferecerem uma abordagem robusta para a avaliação de diferentes cenários, considerando as incertezas inerentes aos parâmetros materiais, geométricos e de carregamento. Conforme apontado por \textcite{pereira_junior_reliability_2023}, essa abordagem permite analisar o desempenho estrutural em distintos níveis de degradação, proporcionando resultados mais realistas e abrangentes.

Contudo, o simples emprego de regressores supervisionados sobre bases sintéticas apresenta uma limitação conceitual que merece destaque. Modelos de regressão convencionais fornecem uma predição \textit{pontual} do tempo de vida útil, ao passo que a decisão de manutenção sob incerteza requer a \textit{distribuição de probabilidade} desse tempo, da qual se extraem quantis inferiores e métricas de risco de cauda. Essa distinção é essencial: a incerteza reportada por um regressor descreve o erro do próprio modelo, e não a variabilidade intrínseca do processo de degradação. Para superar essa limitação, esta dissertação incorpora o arcabouço dos emuladores estocásticos, no qual a distribuição condicional completa da função de estado limite é aproximada por uma Distribuição Lambda Generalizada, cujos parâmetros são representados por Expansão em Caos Polinomial das variáveis de entrada e interpolados continuamente no tempo por um modelo de aprendizado de máquina \cite{ZhuSudret2020, ZhuSudret2021}. Uma vez treinado, o emulador fornece instantaneamente a densidade de probabilidade da margem de segurança em qualquer instante do horizonte de análise, da qual derivam a probabilidade de falha, a distribuição do tempo até a falha, a vida útil remanescente probabilística e as métricas de risco associadas.

O modelo desenvolvido como produto final desta pesquisa consistirá em uma ferramenta computacional, fundamentada em um referencial probabilístico, voltada a subsidiar decisões de manutenção e reforço de estruturas de concreto armado afetadas pela carbonatação, oferecendo suporte quantitativo à gestão da durabilidade estrutural.


\section{Objetivos}

\subsection{Objetivo Geral}

Desenvolver uma metodologia integrada, baseada em simulação de Monte Carlo, emulação estocástica e algoritmos de aprendizado de máquina, para avaliar a confiabilidade de estruturas de concreto armado submetidas à corrosão induzida por carbonatação e, a partir disso, determinar de forma probabilística sua vida útil remanescente (VUR), com o propósito de fornecer uma ferramenta prática, precisa e escalável para análise de risco e suporte à tomada de decisão em manutenção e gestão estrutural.

\subsection{Objetivos Específicos}

\begin{alineas}
    \item Representar numericamente os processos de degradação das propriedades mecânicas do concreto e das armaduras ao longo do tempo, considerando os efeitos da corrosão induzida por carbonatação.
    \item Gerar um banco de dados sintético por meio de simulações de Monte Carlo, incorporando a variabilidade de parâmetros geométricos, mecânicos e de carregamento.
    \item Aplicar algoritmos de aprendizado de máquina para modelar a evolução do índice de confiabilidade ($\beta$) e prever a vida útil remanescente (VUR) das estruturas em função de suas características e condições de exposição.
    \item Desenvolver um emulador estocástico dependente do tempo, capaz de aproximar a distribuição condicional completa da função de estado limite por meio da Distribuição Lambda Generalizada acoplada à Expansão em Caos Polinomial, com interpolação temporal contínua dos parâmetros da distribuição por aprendizado de máquina.
    \item Verificar o emulador proposto em problema de referência com solução de Monte Carlo direta, quantificando a aderência distribucional por testes estatísticos, comparação de quantis e ganho de eficiência computacional.
    \item Derivar da distribuição emulada as grandezas de interesse para a gestão de ativos: a probabilidade de falha dependente do tempo, a distribuição do tempo até a falha, a VUR probabilística e métricas de risco de cauda (VaR e CVaR).
    \item Incorporar métricas de incerteza, como intervalos de predição e largura de banda de confiança, para avaliar a robustez e a confiabilidade dos modelos desenvolvidos.
    \item Validar o desempenho dos modelos preditivos por meio de métricas estatísticas de erro, análises de sensibilidade e simulações independentes, garantindo sua aplicabilidade prática em diferentes cenários estruturais.
\end{alineas}


\section{Justificativa do Tema}

O concreto armado constitui o principal material estrutural utilizado na construção civil brasileira, sendo empregado em praticamente todos os tipos de obras, desde edifícios residenciais e comerciais até pontes, viadutos e barragens. Dessa forma, é possível afirmar que a grande maioria das construções urbanas no país utiliza algum tipo de estrutura de concreto armado, dada sua resistência mecânica, durabilidade e adaptabilidade a diferentes projetos. Essa predominância evidencia não apenas a importância estratégica dessas estruturas para a economia e a sociedade, mas também reforça a necessidade de estudos aprofundados sobre durabilidade, manutenção e vida útil, de modo a assegurar a segurança e a funcionalidade de um vasto parque construído ao longo das últimas décadas.

Por estarem constantemente expostas a ambientes agressivos, essas estruturas tornam-se suscetíveis a processos de degradação que comprometem seu desempenho ao longo do tempo. Dentre esses mecanismos, destaca-se a corrosão das armaduras induzida pela carbonatação, que reduz gradualmente a capacidade resistente das seções estruturais e, consequentemente, a vida útil da estrutura. Considerando a complexidade intrínseca desse fenômeno e a variabilidade dos parâmetros envolvidos, como propriedades dos materiais, condições ambientais e cargas aplicadas, torna-se essencial adotar ferramentas analíticas capazes de representar adequadamente tais incertezas. As abordagens determinísticas tradicionais, embora amplamente utilizadas, não conseguem capturar com precisão o comportamento probabilístico do sistema estrutural frente à degradação progressiva.

Nesse contexto, a aplicação de métodos probabilísticos, como a simulação de Monte Carlo, permite uma avaliação mais realista do tempo de vida útil remanescente das estruturas de concreto armado, incorporando a aleatoriedade dos parâmetros físicos e mecânicos. Além disso, a utilização de algoritmos de aprendizado de máquina possibilita a construção de modelos preditivos eficientes, capazes de generalizar padrões complexos a partir de grandes volumes de dados simulados.

A predição probabilística da vida útil remanescente oferece benefícios diretos significativos: possibilita o planejamento de manutenções baseadas em condições reais da estrutura, promove a alocação otimizada de recursos, aumenta a segurança ao identificar antecipadamente potenciais falhas e contribui para a sustentabilidade, ao permitir a extensão da vida útil das edificações sem comprometer sua integridade estrutural.

Dessa forma, a implementação de uma metodologia que integre simulação estocástica e aprendizado de máquina apresenta-se como uma solução robusta para a estimativa da vida útil de estruturas de concreto armado expostas à corrosão por carbonatação, fornecendo suporte quantitativo para decisões estratégicas de manutenção, reabilitação e gestão de ativos, com impactos diretos na durabilidade, segurança e eficiência do parque construído.

\section{Estrutura do Trabalho}

Além desta introdução, o texto está organizado em quatro capítulos.

O \autoref{cap:referencial} reúne o referencial teórico que sustenta a pesquisa, iniciando pelos fundamentos de confiabilidade estrutural e pelo método de amostragem de Monte Carlo com Hipercubo Latino. Em seguida, apresenta os modelos de resistência à flexão de vigas de concreto armado, com e sem a influência da corrosão das armaduras, e o modelo de avanço da frente de carbonatação. O capítulo trata ainda do conceito de vida útil remanescente e dos fundamentos de aprendizado de máquina, encerrando-se com a apresentação do arcabouço de emuladores estocásticos, que constitui o núcleo metodológico da segunda frente do trabalho.

O \autoref{met} descreve a metodologia adotada, organizada em duas frentes complementares: a construção do banco de dados sintético e o treinamento de modelos de regressão supervisionada, de um lado; e a construção do emulador estocástico dependente do tempo e a derivação, a partir da distribuição emulada, das grandezas de interesse para a gestão de ativos, de outro.

O \autoref{cap:resultados} apresenta os resultados obtidos até o presente estágio da pesquisa, incluindo a comparação entre algoritmos de regressão, a extração das estatísticas de vida útil remanescente, a verificação do emulador estocástico em problema de referência e o estado de desenvolvimento de sua aplicação ao problema da viga sob carbonatação. O capítulo encerra-se com o cronograma revisado de execução da pesquisa.

O \autoref{cap:consideracoes} consolida as considerações parciais que podem ser extraídas do estágio atual do trabalho, registra as limitações reconhecidas e delimita as etapas que restam para a conclusão da dissertação.
